\documentclass[journal]{IEEEtran}

\usepackage[utf8]{inputenc}
\usepackage[T1]{fontenc}
\usepackage{amsmath,amssymb,amsfonts}
\usepackage{graphicx}
\usepackage{booktabs}
\usepackage{array}
\usepackage{url}
\usepackage[hidelinks]{hyperref}
\usepackage{orcidlink}
\usepackage{pgfplots}
\usepackage{tikz}
\usetikzlibrary{arrows.meta,positioning}
\pgfplotsset{compat=1.18}
\graphicspath{{../figures/}}

% IEEEtran hard-codes an em-dash after the abstract and index-terms names; the CAOS
% content standard (ADR-0067) excludes em-dashes, so the separator is a full stop.
\makeatletter
\def\abstract{\normalfont
    \if@twocolumn
      \@IEEEabskeysecsize\bfseries\textit{\abstractname}.\ \relax
    \else
      \bgroup\par\addvspace{0.5\baselineskip}\centering\vspace{-1.78ex}\@IEEEabskeysecsize\textbf{\abstractname}\par\addvspace{0.5\baselineskip}\egroup\quotation\@IEEEabskeysecsize
    \fi\@IEEEgobbleleadPARNLSP}
\def\IEEEkeywords{\normalfont
    \if@twocolumn
      \@IEEEabskeysecsize\bfseries\textit{\IEEEkeywordsname}.\ \relax
    \else
      \bgroup\par\addvspace{0.5\baselineskip}\centering\@IEEEabskeysecsize\textbf{\IEEEkeywordsname}\par\addvspace{0.5\baselineskip}\egroup\quotation\@IEEEabskeysecsize
    \fi\@IEEEgobbleleadPARNLSP}
\makeatother

% Page-1 header block (conventions/manuscript-header-standard.md). The four document
% macros are set by the publishing tool when the version DOI is reserved.
\newcommand{\doctype}{Preprint}
\newcommand{\docversion}{v0.01}
\newcommand{\docdate}{24 September 2026}
\newcommand{\versiondoi}{10.5281/zenodo.NNNNNNNN}
\newcommand{\conceptdoi}{10.5281/zenodo.NNNNNNNN}
\newcommand{\docrepo}{github.com/fsantibanezleal/CAOS\_Enunciado}
\newcommand{\headerblock}{%
\begin{center}
\fbox{\parbox{0.94\textwidth}{\small\raggedright
\textbf{\doctype} \quad$\cdot$\quad \docversion, \docdate \quad$\cdot$\quad Licence CC BY 4.0\\[2pt]
CAOS open-research programme \quad$\cdot$\quad Enunciado, narrative-to-formal measurement\\[2pt]
Version DOI \href{https://doi.org/\versiondoi}{\texttt{\versiondoi}} \quad$\cdot$\quad
concept DOI (always latest) \href{https://doi.org/\conceptdoi}{\texttt{\conceptdoi}}\\[2pt]
Not peer reviewed. Source code, data and computational records: \href{https://\docrepo}{\texttt{\docrepo}}
}}
\end{center}}

% Generated by make_numbers.py from data/artifacts. Do not edit.
\newcommand{\callFailedSentence}{One call failed at the provider after reaching it, and is counted against the model.}
\newcommand{\capAtHighDeepseek}{0/20}
\newcommand{\capAtHighGlmFive}{0/20}
\newcommand{\capAtLowDeepseek}{17/20}
\newcommand{\capAtLowGlmFive}{8/20}
\newcommand{\capCostHighDeepseek}{0.98}
\newcommand{\capCostHighGlmFive}{0.80}
\newcommand{\capCostLowDeepseek}{0.66}
\newcommand{\capCostLowGlmFive}{0.64}
\newcommand{\capFaithfulHighDeepseek}{11/20}
\newcommand{\capFaithfulHighGlmFive}{8/19}
\newcommand{\capFaithfulLowDeepseek}{2/20}
\newcommand{\capFaithfulLowGlmFive}{7/20}
\newcommand{\capHigh}{32768}
\newcommand{\costUSD}{4.01}
\newcommand{\faithfulDeepseek}{0.133}
\newcommand{\faithfulGlmFive}{0.350}
\newcommand{\faithfulGlmFlash}{0.050}
\newcommand{\faithfulHaiku}{0.250}
\newcommand{\faithfulHighDeepseek}{0.297}
\newcommand{\faithfulHighGlmFive}{0.567}
\newcommand{\faithfulHighGlmFlash}{0.236}
\newcommand{\faithfulHighHaiku}{0.402}
\newcommand{\faithfulHighPhi}{0.469}
\newcommand{\faithfulHighQwenFourteen}{0.301}
\newcommand{\faithfulHighSonnet}{0.671}
\newcommand{\faithfulLowDeepseek}{0.053}
\newcommand{\faithfulLowGlmFive}{0.181}
\newcommand{\faithfulLowGlmFlash}{0.009}
\newcommand{\faithfulLowHaiku}{0.142}
\newcommand{\faithfulLowPhi}{0.112}
\newcommand{\faithfulLowQwenFourteen}{0.028}
\newcommand{\faithfulLowSonnet}{0.375}
\newcommand{\faithfulNDeepseek}{4/30}
\newcommand{\faithfulNGlmFive}{7/20}
\newcommand{\faithfulNGlmFlash}{1/20}
\newcommand{\faithfulNHaiku}{10/40}
\newcommand{\faithfulNPhi}{5/20}
\newcommand{\faithfulNQwenFourteen}{2/20}
\newcommand{\faithfulNSonnet}{21/40}
\newcommand{\faithfulPhi}{0.250}
\newcommand{\faithfulQwenFourteen}{0.100}
\newcommand{\faithfulSonnet}{0.525}
\newcommand{\gapAllowedHaiku}{0.000}
\newcommand{\gapAllowedPhi}{+0.050}
\newcommand{\gapAllowedSonnet}{0.000}
\newcommand{\gapDeepseek}{0.000}
\newcommand{\gapGlmFive}{0.000}
\newcommand{\gapGlmFlash}{0.000}
\newcommand{\gapHaiku}{+0.075}
\newcommand{\gapPhi}{+0.100}
\newcommand{\gapQwenFourteen}{+0.100}
\newcommand{\gapSonnet}{+0.050}
\newcommand{\measuredFrom}{2026-09-22}
\newcommand{\measuredTo}{2026-09-24}
\newcommand{\nAtCapAll}{79}
\newcommand{\nCallFailed}{1}
\newcommand{\nCalls}{370}
\newcommand{\nCases}{20}
\newcommand{\nFabricated}{104}
\newcommand{\nFaithfulAll}{53}
\newcommand{\nFaithfulPropertyOnly}{47}
\newcommand{\nGapDefined}{10}
\newcommand{\nGapPositive}{6}
\newcommand{\nHosted}{5}
\newcommand{\nLocal}{11}
\newcommand{\nMeasuredAll}{369}
\newcommand{\nModels}{16}
\newcommand{\nNoAnswerCap}{64}
\newcommand{\nNoUnit}{42}
\newcommand{\nNothingRan}{6}
\newcommand{\nPerModel}{40}
\newcommand{\nProviders}{4}
\newcommand{\nRanAll}{67}
\newcommand{\nRefutedOther}{7}
\newcommand{\nRefutedWhole}{6}
\newcommand{\nRepeatCopies}{0}
\newcommand{\nRepeatIdentical}{0}
\newcommand{\nRepeatPairs}{50}
\newcommand{\nRepeatPairsHosted}{50}
\newcommand{\nRepeatSameClass}{28}
\newcommand{\nRepeatSameFaithful}{39}
\newcommand{\nRepeatSameFaithfulHosted}{39}
\newcommand{\nRepeats}{2}
\newcommand{\nStructFail}{13}
\newcommand{\nStructPass}{6}
\newcommand{\nStructUndecided}{48}
\newcommand{\nTiers}{5}
\newcommand{\nTruncated}{21}
\newcommand{\nUnmeasured}{1}
\newcommand{\nWholeCases}{3}
\newcommand{\nWholeHaiku}{3}
\newcommand{\nWholePhi}{1}
\newcommand{\nWholeReadings}{6}
\newcommand{\nWholeSonnet}{2}
\newcommand{\protocolCap}{8192}
\newcommand{\ranDeepseek}{0.133}
\newcommand{\ranGlmFive}{0.350}
\newcommand{\ranGlmFlash}{0.050}
\newcommand{\ranHaiku}{0.325}
\newcommand{\ranPhi}{0.350}
\newcommand{\ranQwenFourteen}{0.200}
\newcommand{\ranSonnet}{0.575}
\newcommand{\repeatCopiesSentence}{no model did so.}
\newcommand{\wholeCaseList}{opt-006, opt-010, opt-012}
\newcommand{\wilsonWidth}{0.30}

\def\rateticks{16,15,14,13,12,11,10,9,8,7,6,5,4,3,2,1}
\def\ratelabels{{claude-sonnet-5},{claude-haiku-4-5},{glm-5.3},{glm-4.5-flash},{deepseek-v4-pro},{phi4:latest},{qwen3:14b},{deepseek-r1:8b},{qwen2.5-coder:7b},{gemma3:12b},{gemma3:4b},{llama3.1:8b},{mistral:7b},{phi4-mini:latest},{qwen3:4b},{qwen3:8b}} % voice-ok: a model under measurement


\newcommand{\model}[1]{\texttt{#1}}
% A generated table body is read with the primitive input: LaTeX's \input is not expandable,
% and inside a tabular it breaks the \noalign that the booktabs rules need.
\makeatletter\newcommand{\tablebody}[1]{\@@input #1 }\makeatother

\begin{document}

\title{It Ran Is Not It Was Right: The Formalization Gap of \nModels\ Language Models
on Authored Optimization Problems, and What Its Refutations Rest On}

\author{%
  Felipe Santib\'a\~nez-Leal~\orcidlink{0000-0002-0150-3246},~\IEEEmembership{Independent Researcher}%
  \thanks{F.\ Santib\'a\~nez-Leal is with the CAOS open-research programme,
    Santiago, Chile (e-mail: \texttt{fsantibanezleal@ug.uchile.cl};
    ORCID: 0000-0002-0150-3246).}%
  \thanks{The corpus, the per-call records and the derived tables are released at
    \url{https://github.com/fsantibanezleal/CAOS_Enunciado}; the instrument is released as
    the Python packages \texttt{planteo} and \texttt{copela} (MIT); an interactive viewer is
    available at \url{https://enunciado.fasl-work.com}.}%
}

\markboth{\doctype\ \docversion, 2026 \textperiodcentered\ CAOS open-research programme}{Santib\'a\~nez-Leal:
The Formalization Gap on Authored Optimization Problems}

\IEEEaftertitletext{\vspace{-0.6\baselineskip}\headerblock\vspace{0.2\baselineskip}}

\IEEEtitleabstractindextext{%
\begin{abstract}
Language models that turn a problem statement into an optimization model are evaluated almost
everywhere by whether the model runs and reaches the reference optimum. That criterion accepts a
formalization that runs and describes a different problem. We measure the distance between the two
with an instrument whose layers are reported separately and never merged: an executable layer
(the typed document validates and the solver reaches an optimum), a structural layer that proves
equivalence by equal canonical forms and refutes it by optima that differ in the minimising sense,
and a property layer of metamorphic relations. The gap between the rate at which formalizations ran
and the rate at which they also survived both strong layers is the quantity reported. On an authored
corpus of \nCases\ cases in \nTiers\ tiers, with verified references, we ran \nModels\ models from
\nProviders\ providers (\nHosted\ hosted, \nLocal\ open-weight models on one 8~GB GPU), \nRepeats\
repeats per case, \nCalls\ calls. Three results qualify the headline rates. First, the output cap
decides the reasoning models: at \protocolCap\ tokens DeepSeek-V4-Pro was faithful on
\capFaithfulLowDeepseek\ of the same passes and at \capHigh\ on \capFaithfulHighDeepseek. Second, a
refutation inherits the reference's choices: \nWholeReadings\ refutations land exactly on the
reference's optimum with its decisions made integer, in statements that never say whether the
decisions are whole numbers. Third, a model repeats its own outcome on a case only part of the time.
The structural layer decided \nStructPass\ plus \nStructFail\ of the \nRanAll\ candidates that ran,
so most faithful outcomes rest on the property layer alone.
\end{abstract}

\begin{IEEEkeywords}
Autoformalization, optimization modelling, large language models, evaluation, metamorphic testing,
equivalence checking, benchmark validity.
\end{IEEEkeywords}}

\maketitle

\section{Introduction}

\IEEEPARstart{A}{} model that compiles is not necessarily the model that was asked for. In
natural-language-to-Lean statement formalization, where this has been measured with care, every one
of eight systems showed a nonzero gap between compiling and being faithful to the statement, from
3.0 to 29.0 percentage points over 400 graduate-level statements, and the system that compiled most
often (89.5\%) was faithful on 60.5\% \cite{lean2026}. The same study concludes that language-model
judging is a calibrated, conservative aggregate, not an equivalence oracle.

Optimization modelling has the same structure and a weaker practice. Its benchmarks score a
formalization by whether it runs and reaches the reference optimum \cite{survey2025,nl4opt2023}.
Two objections apply. A matching optimum does not establish equivalence, because compensating errors
reach the right number; and the benchmarks themselves carry errors: a recent survey had eleven
experts re-check seven of them and found error rates from at least 8.13\% to at least 54.0\%
\cite{survey2025}. Graph-isomorphism tests over the instance graph of a linear or mixed-integer
program give a stronger equivalence check for that class \cite{orgeval2025}, but they do not reach
formalizations outside it.

We measure the gap between running and being faithful for optimization, on a corpus written for the
purpose, across many models, with oracles that are not language models. The contributions are:
\begin{itemize}
\item an instrument, released as two packages, whose executable, structural and property layers are
reported separately, with the gap between them as the output (Section~\ref{sec:instrument});
\item an authored corpus of \nCases\ cases whose references were solved and checked against their
claimed optima before any model was run (Section~\ref{sec:corpus});
\item a measurement of \nModels\ models, hosted and local, at \nRepeats\ repeats
(Sections~\ref{sec:protocol} and~\ref{sec:results});
\item three results that qualify what such a measurement can claim: the output cap decides the
reasoning models' rates, refutations can rest on choices the statement left open, and a model's
outcome on a case does not always repeat (Section~\ref{sec:results}).
\end{itemize}

\section{Related work}

Autoformalization evaluation has moved from compilation to faithfulness \cite{lean2026,common2025}.
In optimization modelling, benchmarks such as NL4Opt \cite{nl4opt2023} and the sets collected by
\cite{survey2025} use execution accuracy. ORGEval represents a linear or mixed-integer instance as a
bipartite graph and decides equivalence by a Weisfeiler-Lehman test under a sufficient condition
\cite{orgeval2025,shervashidze2011}; the refinement cannot separate all non-isomorphic graphs in
general \cite{cfi1992}. Test oracles and metamorphic relations are established tools where no
reference output is available \cite{barr2015,segura2016}; we use them as the property layer.
Experiment design and simulation modelling are being benchmarked with language-model judges
\cite{scope2026,beams2026}; our layers are deterministic.

\section{The instrument}\label{sec:instrument}

\subsection{A typed representation}

A formalization is a document with quantities, relations and objectives. Each quantity has a role
(decision variable, parameter or derived), a physical dimension, a domain, bounds, and a span that
cites the words of the statement it comes from. Validation rejects a document whose relations mix
dimensions, whose names are used but never declared, whose derived quantities are never defined, or
whose spans quote words the statement does not contain. A valid document emits to Pyomo
\cite{pyomo} and is solved with HiGHS \cite{highs}.

\subsection{Three layers, never merged}

The \emph{executable} layer passes a candidate that validates and whose solve reaches an optimum
(or, for a model with no objective, a feasible point). An unbounded model does not pass.

The \emph{structural} layer compares the candidate $P$ with the case's reference $R$. Equal canonical
forms prove equivalence and pass. Otherwise both are solved; with $s(P)=+1$ for minimisation and
$-1$ for maximisation, and $z^*$ the optimal value,
\begin{equation}
s(P)\,z^*(P) \neq s(R)\,z^*(R) \;\Rightarrow\; P \text{ and } R \text{ are different models,}
\end{equation}
so a differing optimum refutes. A matching optimum decides nothing, because compensating errors reach
the right number; the layer then reports undecided.

The \emph{property} layer applies metamorphic relations that any correct formalization of the class
satisfies \cite{segura2016}: scaling the objective does not move the argmin, a redundant constraint
does not change the feasible set, tightening a constraint cannot improve the optimum, relaxing one
cannot worsen it, and permuting rows and columns changes nothing. A violated relation refutes.

A candidate is \emph{faithful} when it ran, neither strong layer refuted it, and at least one of them
passed it. With $R_{\mathrm{ran}}(m)$ and $R_{\mathrm{faithful}}(m)$ the rates of model $m$ over its
calls, the reported quantity is the gap
\begin{equation}
\Delta_m = R_{\mathrm{ran}}(m) - R_{\mathrm{faithful}}(m),
\end{equation}
with each rate carried by a 95\% Wilson interval \cite{wilson1927,agresti1998},
\begin{equation}
\frac{\hat p + \frac{z^2}{2n} \pm z\sqrt{\frac{\hat p(1-\hat p)}{n} + \frac{z^2}{4n^2}}}{1 + \frac{z^2}{n}}, \qquad z = 1.96 .
\end{equation}
A candidate the solver cannot express leaves both denominators and is reported as unmeasured. A
language-model judge is recorded for comparability with published work and never counted.

\section{The corpus}\label{sec:corpus}

The \nCases\ cases are authored, four per tier: direct statements (the controls), composed ones
(a unit conversion, a derived quantity, a distractor), structured ones (coupled decisions, balances
across periods, ratios), discrete ones (integrality, fixed charges, minimum runs, assignments) and
underspecified ones, where the statement does not determine the model. Each case names the trap it
was built around. Before any model was run, every reference was validated and solved, every claimed
optimum was compared with the solver's (three of twenty claims were wrong and were corrected), and
every metamorphic relation was checked on the reference. The corpus is authored rather than drawn
from the public sets because of their measured error rates \cite{survey2025} and their licences,
several of which forbid redistribution.

\section{Protocol}\label{sec:protocol}

Every model received the same statement, the same schema and the same instruction, at an output cap
of \protocolCap\ tokens, temperature 0 or the provider's greedy switch where one exists, a fixed
seed, and high reasoning effort where a model accepts the setting. Local models ran through Ollama
in a context sized to hold the prompt and the whole cap, with reasoning switched off where the
model's template honours the switch. Each case was run \nRepeats\ times per model, \nCalls\ calls in
total between \measuredFrom\ and \measuredTo, at a list-price cost of \costUSD\ USD. Every call is
appended to a per-call record with the model version, fingerprint, prompt digest, the outcome of
every layer, the scoring software's version and, for later calls, the parsed document. A call that
never reached the provider (a refused connection or rejected credentials) stops the run and is not
recorded; a provider's error response is recorded as a failed call. The two reasoning models were
also run at a \capHigh-token cap, in a separate record so that neither cap could overwrite the other.

\section{Results}\label{sec:results}

\subsection{The two rates and the gap}

Table~\ref{tab:models} and Fig.~\ref{fig:rates} give each model's rates. Over all models,
\nRanAll\ of \nMeasuredAll\ measured candidates ran and \nFaithfulAll\ were faithful; \nNothingRan\
models produced no candidate that ran, so their gap is undefined. Of the \nGapDefined\ models with a
defined gap, \nGapPositive\ are positive. At $n = \nPerModel$ per model a rate of 0.5 carries an
interval about \wilsonWidth\ wide, so most pairs of models cannot be ranked.

\begin{table*}[t]
\centering
\caption{Rates per model, 95\% Wilson intervals, over \nPerModel\ calls per model.
``At cap'' counts calls that billed exactly the output cap.}
\label{tab:models}
\begin{tabular}{llcccc}
\toprule
Model & Provider & Ran & Faithful & Gap & At cap \\
\midrule
\tablebody{table-models.tex}
\bottomrule
\end{tabular}
\end{table*}

\begin{figure}[t]
\centering
\begin{tikzpicture}
\begin{axis}[
  width=\columnwidth, height=9.5cm,
  xmin=0, xmax=1, xlabel={rate},
  ytick=\rateticks, yticklabels/.expanded=\ratelabels,
  yticklabel style={font=\scriptsize\ttfamily},
  xticklabel style={font=\scriptsize}, xlabel style={font=\small},
  ymin=0.3, ymax=\nModels+0.7, grid=major, grid style={gray!20},
  legend style={font=\scriptsize, at={(0.98,0.98)}, anchor=north east},
]
\addplot+[only marks, mark=*, mark size=1.6pt, black,
  error bars/.cd, x dir=both, x explicit]
  table[x=ran, y expr=\thisrow{i}+0.15, x error minus expr=\thisrow{ran}-\thisrow{ranlo},
        x error plus expr=\thisrow{ranhi}-\thisrow{ran}]{../figures/rates.dat};
\addlegendentry{ran}
\addplot+[only marks, mark=square*, mark size=1.6pt, blue!70!black,
  error bars/.cd, x dir=both, x explicit]
  table[x=faithful, y expr=\thisrow{i}-0.15, x error minus expr=\thisrow{faithful}-\thisrow{faithfullo},
        x error plus expr=\thisrow{faithfulhi}-\thisrow{faithful}]{../figures/rates.dat};
\addlegendentry{faithful}
\end{axis}
\end{tikzpicture}
\caption{Ran and faithful rates per model with 95\% Wilson intervals; the horizontal distance
between the two marks of a model is its gap.}
\label{fig:rates}
\end{figure}

\subsection{Where the calls fail}

The dominant failure of the local models is fabricated provenance: a span that cites words the
statement does not contain, \nFabricated\ calls in all. The dominant failure of the reasoning models
is the cap: \nNoAnswerCap\ calls spent all \protocolCap\ tokens reasoning and returned no document,
and \nTruncated\ more were cut off mid-document. \nNoUnit\ calls left a constant without its unit.
\callFailedSentence

\subsection{The cap decides the reasoning models}

Table~\ref{tab:cap} compares the two reasoning models at \protocolCap\ and \capHigh\ tokens over the
same passes. DeepSeek-V4-Pro was faithful on \capFaithfulLowDeepseek\ at the protocol's cap, with
\capAtLowDeepseek\ calls reaching it, and on \capFaithfulHighDeepseek\ at \capHigh, with none
reaching it. At the lower cap its failures are about the cap, not about how it formalizes, and a
single-cap table would rank it among the weakest models for that reason.

\begin{table}[t]
\centering
\caption{The same passes at two output caps.}
\label{tab:cap}
\begin{tabular}{lrrrrr}
\toprule
Model & Cap & Ran & Faithful & At cap & USD \\
\midrule
\tablebody{table-cap.tex}
\bottomrule
\end{tabular}
\end{table}

\subsection{What the refutations rest on}

The structural layer decided \nStructPass\ candidates by equal canonical forms and \nStructFail\ by
refutation, of the \nRanAll\ that ran; the other \nStructUndecided\ were undecided, so
\nFaithfulPropertyOnly\ of the \nFaithfulAll\ faithful outcomes rest on the property layer alone.

A refutation says the candidate and the reference are different models. It does not say the
candidate misread the statement, because the reference also chose. Re-solving every reference with
its real decision variables made integer moves the optimum in \nWholeCases\ cases (\wholeCaseList).
Two of them ask for the most night shifts a wage bill allows and for the best margin on pumps and
valves; neither statement says whether a shift or a pump comes in whole units, and both references
are continuous. \nWholeReadings\ refutations land exactly on the reference's whole-number optimum.
Counting them as readings the statement allows, and faithful where the property layer passed them,
moves Claude Sonnet~5's gap from \gapSonnet\ to \gapAllowedSonnet\ and Claude Haiku~4.5's from % voice-ok: Claude names models under measurement
\gapHaiku\ to \gapAllowedHaiku. A matching optimum does not prove the candidate is the integer model;
the refutation remains in the published rates, and the reading is reported beside it.

\subsection{Run to run}

Table~\ref{tab:runs} compares each model's second pass over a case with its first. Of \nRepeatPairs\
repeated cases, \nRepeatSameFaithful\ reached the same faithful outcome and \nRepeatSameClass\ the
same failure class; \nRepeatIdentical\ returned the same response byte for byte. Among the hosted
models, \nRepeatSameFaithfulHosted\ of \nRepeatPairsHosted\ repeated the faithful outcome. A model
whose every repeat returns its first response contributes one sample counted twice;
\repeatCopiesSentence

\begin{table}[t]
\centering
\caption{Later repeats of a case compared with the first: identical response, same failure class,
same faithful outcome.}
\label{tab:runs}
\begin{tabular}{lrrrr}
\toprule
Model & Pairs & Identical & Class & Faithful \\
\midrule
\tablebody{table-runs.tex}
\bottomrule
\end{tabular}
\end{table}

\section{Limitations}

The statements, the references and the checks were written by one author; a statement and its
reference share their reading, and a second author formalizing the cases blind would be needed to
measure how often a reference is one reading among several. The integer results above are two
instances of that. Twenty cases at two repeats cannot rank most pairs of models, and repeats of a
case are not independent samples. The structural layer compares the typed document; a stronger
canonical form over the linear rows exists and was not applied to the recorded candidates. The
per-call records written before the scoring software began recording its own version were scored by
three versions of it, whose differences change no recorded outcome but one unbounded candidate,
which is classified separately. The measurement covers linear and mixed-integer optimization only.

\section{Conclusion}

Execution accuracy overstates what a formalization gets right, and so can a stronger check applied
without asking what its reference chose. Reported separately, with the cap, the refutations' basis
and the repeats stated beside the rates, the gap between running and being faithful is measurable
on a small authored corpus; ranking models by it needs more cases and more repeats than any single
pass provides.

\section*{Data and code availability}

The corpus, the per-call records, the derived tables and the pipeline that produces them are at
\url{https://github.com/fsantibanezleal/CAOS_Enunciado}; every number in this paper is transcribed by
a script from those tables. The instrument is released as \texttt{planteo} and \texttt{copela} on
PyPI. Reproducing a table requires no model call; reproducing a call requires the provider and is
not deterministic for hosted models.

\bibliographystyle{IEEEtran}
\bibliography{refs}

\end{document}
